% !TEX root = IE3_expJ_report.tex
\documentclass[lualatex,a4paper,10pt]{jlreq}
% Shared LuaLaTeX preamble.
\usepackage{luatexja-fontspec}
\setmainfont{TeX Gyre Termes}
\setsansfont{TeX Gyre Heros}
\setmonofont{Latin Modern Mono}
\setmainjfont[BoldFont=HaranoAjiGothic-Medium]{HaranoAjiMincho-Regular}
\setsansjfont[BoldFont=HaranoAjiGothic-Bold]{HaranoAjiGothic-Medium}
\usepackage[top=22mm,bottom=22mm,left=23mm,right=23mm]{geometry}
\usepackage{amsmath,amssymb,booktabs,array,longtable,tabularx}
\usepackage{graphicx,xcolor,tikz,circuitikz}
\usetikzlibrary{arrows.meta,positioning,calc}
\usepackage{enumitem,fancyhdr,listings,needspace}
\usepackage[unicode,colorlinks=true,linkcolor=labblue,urlcolor=labteal]{hyperref}
\usepackage{bookmark}
\definecolor{labblue}{RGB}{30,70,112}
\definecolor{labteal}{RGB}{0,100,105}
\definecolor{labgray}{gray}{0.95}
\ModifyHeading{section}{font={\large\sffamily\bfseries\color{labblue}}}
\ModifyHeading{subsection}{font={\normalsize\sffamily\bfseries\color{labteal}}}
\setlength{\parindent}{1\zw}
\setlength{\parskip}{3pt}
\setlength{\emergencystretch}{2em}
\renewcommand{\arraystretch}{1.35}
\setlist{itemsep=3pt,topsep=4pt,leftmargin=2\zw}
\newcolumntype{Y}{>{\raggedright\arraybackslash}X}
\newcolumntype{C}[1]{>{\centering\arraybackslash}p{#1}}
\pagestyle{fancy}
\fancyhf{}
\fancyhead[L]{\small\color{labblue} IE3 後期・電子工学実験}
\fancyhead[R]{\small テーマ\LabCode}
\fancyfoot[C]{\thepage}
\setlength{\headheight}{16pt}
\DeclareOldFontCommand{\rm}{\normalfont\rmfamily}{\mathrm}
\lstset{basicstyle=\small\ttfamily,columns=fullflexible,keepspaces=true,
 breaklines=true,frame=single,backgroundcolor=\color{labgray},
 numbers=left,numberstyle=\tiny,showstringspaces=false,aboveskip=8pt,belowskip=8pt}
\newcommand{\blank}{\rule{22mm}{0.3pt}}
\newcommand{\answerline}{\par\noindent\rule{\linewidth}{0.25pt}\par}
\newcommand{\writing}[1]{\par\Needspace{\dimexpr#1+5mm\relax}\noindent%
 \fcolorbox{labblue!35}{white}{\parbox[c][#1][t]{\dimexpr\linewidth-2\fboxsep-2\fboxrule\relax}{\vspace{2mm}\noindent\strut\hfill\par}}\par}
\newcommand{\hint}[1]{\par\smallskip\noindent{\sffamily\bfseries\color{labteal}記述の視点：}#1\par}
\newcommand{\note}[1]{\par\smallskip\noindent\fcolorbox{labblue!45}{labblue!5}{%
 \parbox{\dimexpr\linewidth-2\fboxsep-2\fboxrule\relax}{#1}}\par\smallskip}
\newcommand{\eqerror}{\varepsilon=\frac{x_{\mathrm{meas}}-x_{\mathrm{th}}}{x_{\mathrm{th}}}\times100\ \%}
\newcommand{\LabSource}{徳山工業高等専門学校 情報電子工学科，
 『2026年度 電子工学実験（3年後期）テキスト』，テーマ\LabCode，
 \expandafter\path\expandafter{\SourcePdf}。}
\newcommand{\reporttitle}{
 \noindent{\small\sffamily\color{labteal}実験報告書・記入ガイド}\par
 \noindent{\LARGE\sffamily\bfseries \LabCode\quad\LabTitle}\par\smallskip
 \noindent 氏名：\blank\quad 班：\blank\quad 実験日：\blank\par
 \note{目的・原理・手順を確認し、実測値と自分の考察を記入する。
 考察のガイドは、検討する事象・使う測定結果・記述の方針を示す。}
}
\newcommand{\prelabtitle}{
 \noindent{\small\sffamily\color{labteal}事前学習・前提知識プリント}\par
 \noindent{\LARGE\sffamily\bfseries \LabCode\quad\LabTitle}\par\smallskip
 \noindent 氏名：\blank\quad 班：\blank\quad 予習日：\blank\par
 \note{用語、回路の動き、計算、測定表への記入の順に読む。
 例題の値は説明用であり、実験結果ではない。}
}
\newcommand{\tocrow}[3]{\hyperref[#1]{#2}& #3 &\pageref*{#1}\\[2mm]}
\newcommand{\documentmap}[1]{
 \section*{目次と概要}
 \noindent 青色の項目をクリックすると該当箇所へ移動する。\par
 {\small\begin{longtable}{@{}p{0.29\linewidth}p{0.61\linewidth}r@{}}
 \toprule {\color{labblue}\bfseries 項目} & {\color{labblue}\bfseries 読むと分かること} & 頁\\\midrule
 \endhead #1\bottomrule
 \end{longtable}}
 \clearpage
}
\newenvironment{terms}{\par\smallskip\begin{longtable}{@{}p{0.23\linewidth}p{0.73\linewidth}@{}}
 \toprule {\color{labteal}\bfseries 用語}&{\color{labteal}\bfseries 意味・回路での役割}\\\midrule\endhead}
 {\bottomrule\end{longtable}}
\newcommand{\term}[2]{\textbf{#1}&#2\\[2mm]}
\newcommand{\guidepoint}[4]{
 \Needspace{32mm}\subsection{#1}
 \noindent{\bfseries\color{labteal}考える事象：}#2\par
 \noindent{\bfseries\color{labteal}参照する結果：}#3\par
 \noindent{\bfseries\color{labteal}記述の方針：}#4\par
}
\newcommand{\reportclosing}[1]{
 \clearpage\section{結論のまとめ方}\label{sec:closing}
 \hint{#1}
 \ConclusionGuide
 \begin{enumerate}
 \item 目的に対応する最も重要な実測結果を、測定条件と数値を添えて述べる。
 \item 原理と一致した点・一致しなかった点を示す。原因は確認済みの事実と推測を分ける。
 \item 実施範囲で言えることと、未確認の条件を一文ずつまとめる。
 \end{enumerate}
 \note{書き出しの型：「○○の条件で△△を測定したところ、□□となった。
 これは◇◇と整合する。ただし××は確認しておらず、一般化には追加測定が必要である。」
 空欄は自分の実測値と根拠で埋める。}
 \noindent 結論（実験後に記入）：\writing{30mm}
 \noindent 改善案・次に確認したい条件：\writing{16mm}
 \section*{参考文献}\label{sec:references}
 \begin{enumerate}
 \item \LabSource
 \item 使用したデータシート・書籍・ウェブ資料（著者、題名、該当箇所、URL、参照日）を追加する。
 \end{enumerate}
}
\newcommand{\prelabclosing}{
 \section*{準備・出典}\label{sec:prelabsource}
 \begin{itemize}[nosep]
 \item 資料、筆記具、記録用紙、電卓と、実験条件・機器設定の記録欄。
 \item 確認問題の解答と、分からない用語・回路点への具体的な質問。
 \end{itemize}
 {\small\noindent\LabSource\par}
}
\newcommand{\simpleflow}[1]{
 \begin{center}
 \begin{tikzpicture}[node distance=5mm and 5mm,
 every node/.style={font=\small},box/.style={draw=labblue,fill=labblue!4,rounded corners=1mm,
 align=center,minimum height=10mm,text width=30mm}]
 #1
 \end{tikzpicture}
 \end{center}
}

\newcommand{\LabCode}{J}
\newcommand{\LabTitle}{PLCによるシーケンス制御}
\newcommand{\SourcePdf}{IE3_expJ_A4.pdf}
\begin{document}
\reportcover\reflection
\section{目的}
PLCで接点論理、自己保持、内部リレー、タイマ、カウンタを実装し、ラダー図と入出力の対応を理解する。ベルトコンベヤ等のシーケンスを確認する。
\section{原理}
PLCは入力を取り込み、ラダープログラムを順に実行し、出力を更新する処理を繰り返す。
ラダーのa接点はビットが1なら成立、b接点は0なら成立する。
物理スイッチのNO/NCとプログラムのa/bを区別する。
自己保持では出力や内部リレーの接点を使って過去の操作を保持する。
停止優先の自己保持は
\[
 Q_{\rm next}=(S\lor Q_{\rm prev})\land\lnot R
\]
と表せる。開始と停止が同時に成立する場合も検討する。
\section{装置・割付}
OMRON CP1L-L14DR-D、CX-Programmer、24 V電源、スイッチ、ランプ、ベルトコンベヤを使用する。
\par\noindent\begin{tabularx}{\linewidth}{YYYY}
\toprule 入力 & アドレス & 出力 & アドレス\\\midrule
SW1（a）&0.00&L1&100.00\\SW2（a）&0.01&L2&100.01\\
SW3（a）&0.02&L3&100.02\\SW4（b）&0.03&ベルト&100.03\\\bottomrule
\end{tabularx}
実際の配線図、端子COM、ソフト版、機器番号：
\writing{18mm}
\section{方法}
\begin{enumerate}
\item 電源OFFで配線し、指定の線色（赤24 V、青GND、黄その他）と端子を確認する。一端子に2本を超えて接続しない。
\item OnOffで物理a/bスイッチとラダーa/b接点の4組を試す。LogicでAND、OR、NOTを試す。
\item Parking1/2で3車室の満車表示L3と、一つ以上空きがある表示L1を実現する。
\item 基本の自己保持を確認し、SW1で起動、SW2で停止するSelfHoldingを作る。
\item IntRelayでSW1・SW2の操作履歴を内部リレーに保持し、両方の操作後にベルトを動かす。SW3で全履歴を解除する。
\item TimerでSW1操作から3秒後に動作し、途中のSW2で取消す。CounterでSW1を3回操作すると動作し、SW2でリセットする。
\end{enumerate}
\clearpage
\section{結果}
表1：物理接点とプログラム接点。操作前後の入力ビットと出力を記す。
\par\noindent\begin{tabularx}{\linewidth}{YYYY}
\toprule 物理接点 & ラダー接点 & 非操作時 & 操作時\\\midrule
a&a&&\\a&b&&\\b&a&&\\b&b&&\\\bottomrule
\end{tabularx}
表2：論理と駐車場。1は入力ビットON。L1は少なくとも一室空き、L3は満車。
\par\noindent\begin{tabularx}{\linewidth}{YYYYY}
\toprule SW1 & SW2 & SW3 & L1観測 & L3観測\\\midrule
0&0&0&&\\0&0&1&&\\0&1&0&&\\0&1&1&&\\
1&0&0&&\\1&0&1&&\\1&1&0&&\\1&1&1&&\\\bottomrule
\end{tabularx}
表3：AND/OR/NOTの理論と観測（必要に応じ列を増やす）。
\par\noindent\begin{tabularx}{\linewidth}{YYYY}
\toprule 入力A,B & AND観測 & OR観測 & NOT A観測\\\midrule
0,0&&&\\0,1&&&\\1,0&&&\\1,1&&&\\\bottomrule
\end{tabularx}
\clearpage
表4：シーケンスの境界試験。
\par\noindent\begin{tabularx}{\linewidth}{YYY}
\toprule プログラム・操作 & 期待動作 & 観測結果\\\midrule
自己保持：起動後に手を離す&運転継続&\\
自己保持：停止／同時入力&停止優先を確認&\\
内部リレー：SW2→SW1&両履歴成立後に運転&\\
内部リレー：SW3&運転と履歴解除&\\
タイマ：SW1後3秒&動作&\\
タイマ：3秒未満でSW2&取消し&\\
カウンタ：1、2、3回&3回目に動作&\\
カウンタ：SW2&計数リセット&\\\bottomrule
\end{tabularx}
図1：各プログラムのラダー図とオンラインモニタ画面（別紙追加可）。
\writing{42mm}
図2：自己保持、タイマ、カウンタの入力・出力時系列。
\writing{35mm}
\section{考察・課題}
\begin{enumerate}
\item 自己保持が記憶として働く理由と、内部リレーを用いる意味を説明する。
\item 停止接点の位置を変えた場合の同時入力を比較し、停止優先と開始優先を論理式で示す。
\item TIMの設定\texttt{\#30}（0.1秒単位）とCNTの\texttt{\#3}を説明する。連続ONと3回のエッジを区別する。
\item スキャン周期、実行順序、同一コイルへの重複書込みが動作に与える影響を述べる。
\item 応用課題を実施した場合は信号機などの仕様・ラダー・確認表を追加する。
\end{enumerate}
\writing{22mm}
\reportclosing{接点論理から時間・回数制御までの確認結果をまとめる。}
\end{document}
